% !TEX root = ../../tree_notes.tex
% =====================================================================
%  5.1  First-In / Last-Out Times and the Bracket Sequence   sec:c24-times
%  Exports: def:c24-times  thm:c24-anc  prop:c24-bracket
%           lst:c24-stamp  ex:c24-times-run
%           prob:c24-throne  lst:c24-throne  prob:c24-color (practice)
% =====================================================================
\section{First-In / Last-Out Times and Bracket Sequences}\label{sec:c24-times}

\begin{previouslybox}
\begin{enumerate}[nosep, leftmargin=1.6em]
  \item In a preorder DFS, in what order are $G$'s nodes visited, and why is the
        \emph{parent guard} mandatory on an undirected \code{adj}? (\secref{c04-rooting},
        \thmref{c04-dfs-orient})
  \item What is $\size[v]$ for every node of $G$ rooted at $0$, and why does a leaf have
        $\size=1$? (\figref{c04-running-g}, Trap~2 of \secref{c04-rooting})
  \item Why did we say the \emph{parent array} is order-independent but the timestamps are
        \emph{not}? (\secref{c04-rooting}, \chref{c24} teaser)
  \item State the preorder / postorder pair for the binary running tree $R$. (\chref{c03})
\end{enumerate}
\end{previouslybox}

\begin{topicintro}
A DFS already visits each node twice: once when it \term{arrives} (``first-in'') and once
when it \term{departs} after finishing the whole subtree (``last-out''). If you write down a
counter at those two moments you get two integers per node, $\tin(v)$ and $\tout(v)$, and
they encode \emph{everything about nesting}. Node $u$ is an ancestor of $v$ if and only if
$u$'s life-span $[\tin(u), \tout(u))$ \emph{contains} $v$'s -- so ancestry, the question
that naively costs a walk up the tree, becomes a single $O(1)$ integer comparison. Erase the
labels and write an open bracket at each arrival, a close at each departure, and the whole
tree collapses to a \term{balanced bracket sequence}, the same object the string chapters
encode and decode. This section builds the stamp, the containment test, and the bracket
view; \secref{c24-flatten} then exploits the \emph{order} of $\tin$ to put subtrees on a
line.
\textbf{Prerequisites.} \secref{c04-rooting} (the parent-guarded DFS, $\size$); \chref{c03}
(pre/post order). This section \emph{adds a clock} to the DFS you already wrote -- nothing
else.

\textbf{After this section you can}
\begin{itemize}[nosep]
  \item stamp $\tin$ (entry) and $\tout$ (exit) on every node in one DFS;   \hfill\lv{2}
  \item decide ``is $u$ an ancestor of $v$?'' in $O(1)$ via containment;    \hfill\lv{2}--\lv{3}
  \item read and write the balanced-bracket / contour encoding of a tree;   \hfill\lv{3}
  \item order nodes by birth time and answer ``who precedes whom in the walk''. \hfill\lv{3}
\end{itemize}
\end{topicintro}

\begin{recognitionbox}[Pattern card -- timestamps, ancestry, brackets]
\reckey{the input is a rooted tree and the question is about \emph{nesting}: ``is $u$ an
ancestor of $v$'', ``sort/return nodes by the order a DFS meets them'', ``the whole
subtree of $x$'', or the data \emph{itself} is a string of \emph{balanced parentheses} such
as \code{(())(())}.}
\recuse{
  \item repeated ``ancestor of'' / ``in the subtree of'' yes/no questions -- the $O(1)$
        containment test (\thmref{c24-anc});
  \item a nested-parenthesis string you must parse or reconstruct -- the bracket view below;
  \item ``list nodes in visit order with some nodes filtered out'' -- preorder $=\tin$ order;
  \item as a \emph{prelude} to any subtree range query (that is \secref{c24-flatten}).
}
\recnotuse{
  \item a \emph{single} bottom-up aggregate per node, no queries -- plain DFS from
        \secref{c02-bottomup} needs no timestamps;
  \item you need the range \emph{answered with updates} -- that reduction is
        \secref{c24-flatten}, then the Fenwick/segment chapter;
  \item you need \emph{distances}/LCA -- the $2n{-}1$ variant is \secref{c24-dist}.
}
One-line test: \emph{does the question ask who \textbf{contains} whom, or hand me nested
brackets? Then stamp times -- this section.}
\end{recognitionbox}

\subsection{Intuition}

Run the DFS of \secref{c04-rooting} on $G$, but carry a stopwatch. The instant you
\emph{first} step onto a node you read the clock and record $\tin$; the instant you are about
to \emph{back out} of it for the last time you read the clock again and record $\tout$.
Because a subtree is entered once and left once, and everything inside it is entered and
left strictly in between, $[\tin(u), \tout(u))$ is a closed interval that \emph{either
contains a descendant's whole interval or misses it entirely}. Two intervals of this kind
never partially overlap -- they nest or they're disjoint. That is the entire magic.

\begin{figure}[H]
\centering
\begin{tikzpicture}[diagram, x=13mm, y=9mm]
  % timeline
  \draw[-Stealth, thick] (0,0) -- (8.4,0) node[right, font=\scriptsize] {clock};
  \foreach \t in {0,...,7} \draw (\t,0.12) -- (\t,-0.12) node[below, font=\scriptsize] {\t};
  % intervals for nodes 0,1,3 (nesting) and 2,4,5,6
  \draw[accentorange, line width=1.6pt] (0,1.5) -- (7,1.5) node[midway,above,font=\scriptsize] {$0:[\,0,7)$};
  \draw[primaryblue, line width=1.6pt] (1,1.0) -- (4,1.0) node[midway,above,font=\scriptsize] {$1:[1,4)$};
  \draw[successgreen, line width=1.6pt] (2,0.55) -- (3,0.55) node[above,font=\tiny,xshift=6pt] {$4$};
  \draw[successgreen, line width=1.6pt] (3,0.55) -- (4,0.55) node[above,font=\tiny,xshift=8pt] {$5$};
  \draw[deepblue, line width=1.6pt] (4,1.0) -- (5,1.0) node[above,font=\scriptsize,xshift=3pt] {$2$};
  \draw[primaryblue, line width=1.6pt] (5,1.0) -- (7,1.0) node[above,font=\scriptsize,xshift=3pt] {$3:[5,7)$};
  \draw[successgreen, line width=1.6pt] (6,0.55) -- (7,0.55) node[above,font=\tiny,xshift=6pt] {$6$};
  \node[annot, align=left, font=\scriptsize] at (0.1,2.1)
    {$\subtree(1)=[1,4)$ nests $4:[2,3)$ and $5:[3,4)$;\\
     it is disjoint from $2:[4,5)$ and from $3:[5,7)$.\\
     $1$ is an ancestor of $4$ (contained); of $2$ it is \emph{not}.};
\end{tikzpicture}
\caption{The intervals are the tree. Reading $\tin/\tout$ off the DFS clock for $G$ turns
each node into a segment on a line; segments \emph{nest} exactly when the nodes are
ancestor/descendant and are \emph{disjoint} for cousins. Ancestry becomes interval
containment, checkable with two integer comparisons.}\label{fig:c24-intervals}
\end{figure}

Now \emph{erase the node numbers} and, walking the same DFS, print \code{(} when you enter a
node and \code{)} when you leave. For $G$ you get
$\code{((()())(()())())}$ -- a \term{balanced bracket sequence} where matched pairs are
ancestor/descendant. This is the same nesting you saw in the string encoding
$0(1(4)(5))(2)(3(6))$ of \secref{c04-formats}, only stripped to bare parentheses; the depth
of any position is the number of unmatched opens, i.e.\ the current DFS depth.

\subsection{Definitions and Notation}

\begin{definition}{Entry and exit times; the Euler interval}{c24-times}
Fix a rooted tree and a DFS that visits children in some fixed order. Maintain a counter
$c$, starting at $0$, incremented by one each time a node is \emph{first entered}. The
\term{entry time} $\tin(v)$ is the counter's value when $v$ is first reached (so
$\tin$ is exactly $v$'s \textbf{preorder rank}); the \term{exit time} is
$\tout(v) = \tin(v) + \size[v]$, where $\size[v]$ is $v$'s subtree size. Node $v$'s
\term{Euler interval} is the half-open range $[\tin(v), \tout(v))$. A node with no children
has $\tout(v) = \tin(v) + 1$, an interval of length one.
\end{definition}

\noindent\textbf{In plain words.} $\tin(v)$ is ``the number of nodes the DFS had met before
meeting $v$'', and $\tout(v)$ is ``$\tin$ plus how many nodes live in $v$'s subtree''. The
interval $[\tin(v), \tout(v))$ is precisely the run of preorder positions occupied by $v$ and
all its descendants. We chose the ``$\tout = \tin + \size$'' convention so the interval
length equals the subtree size and the flattened array is exactly $n$ long
(\secref{c24-flatten}); a ``separate clock on exit'' variant that yields times up to $2n$
is mentioned in the variations.

\begin{theorem}{Ancestor test by interval containment}{c24-anc}
For distinct nodes $u, v$ in a rooted tree, $u$ is an ancestor of $v$ if and only if
\[
   \tin(u) \;\le\; \tin(v) \;<\; \tout(u).
\]
Equivalently $\subtree(v) \subseteq \subtree(u)$ iff $[\tin(v), \tout(v)) \subseteq
[\tin(u), \tout(u))$. Two Euler intervals are either nested or disjoint -- never partially
overlapping.
\end{theorem}
\begin{proofbox}
The DFS enters $u$ at time $\tin(u)$ and, because $u$ is entered once and left once after
all of $\subtree(u)$ is processed, occupies the contiguous block $[\tin(u), \tin(u){+}\size[u))$
of the preorder; by definition $\tout(u)=\tin(u)+\size[u]$. If $u$ is an ancestor of $v$,
then $v \in \subtree(u)$, so $v$'s preorder rank lies inside that block: $\tin(u) \le
\tin(v) < \tout(u)$, and moreover $\subtree(v) \subseteq \subtree(u)$ so
$[\tin(v),\tout(v))$ fits wholly inside $[\tin(u),\tout(u))$. Conversely, if $\tin(v) <
\tout(u)$ and $\tin(u) \le \tin(v)$, then $v$'s entry happens while $u$ is still ``open''
(between entering and leaving $u$), which for a DFS means $v$ was reached \emph{through}
$u$, i.e.\ $u$ is on the path from the root to $v$. Partial overlap is impossible because
the DFS's stack discipline means a node's entire interval is written between the moment it
is entered and the moment it is left; a second interval begun before the first ends must
also end before it (nested), and one begun after the first ends is disjoint.
\end{proofbox}
\begin{insightbox}
The proof uses only the ``enter once, leave once, descendants in between'' property -- so
$\tin/\tout$ ancestry is valid for \textbf{any} child order. But the \emph{numeric} values
of $\tin$ depend on the order you visit children; re-order the children and the intervals
reshuffle (while staying valid). ``Order-independent'' is a property of $\parent$
(\secref{c04-rooting}), \emph{not} of $\tin$ -- the asymmetry we flagged earlier.
\end{insightbox}

\begin{proposition}{Balanced-bracket characterisation}{c24-bracket}
Writing \code{(} on entry and \code{)} on exit yields a balanced parenthesis string of
length $2n$ whose matching-pair structure is (isomorphic to) the tree: each node is one
matched pair, and $u$'s pair encloses $v$'s pair iff $u$ is an ancestor of $v$. Conversely
any balanced string of $n$ pairs is the bracket image of some rooted ordered tree.
\end{proposition}
\begin{proofbox}
Every node contributes exactly one \code{(} (at $\tin$) and one \code{)} (at $\tout$), so
$2n$ symbols. At any prefix the number of opens $\ge$ closes (you cannot leave a node you
have not entered), and totals are equal -- the definition of balanced. Nesting of pairs
mirrors interval nesting (\thmref{c24-anc}). For the converse, parse a balanced string with
a stack: an \code{(} pushes a new child, a \code{)} pops; the resulting parent/child links
form a tree whose DFS reproduces the string.
\end{proofbox}

\subsection{The Algorithm: one DFS that stamps \texorpdfstring{$\tin/\tout$}{tin/tout}}

\begin{lstlisting}[caption={Stamp entry/exit times and answer ancestry in O(1)}, label={lst:c24-stamp}]
#include <bits/stdc++.h>
using namespace std;

vector<vector<int>> adj;
vector<int> tin_, tout_, sz_;
int timer_;

// preorder size is enough for tout = tin + size, but the classic form
// stamps BOTH on a single clock that ticks on entry AND exit:
void euler(int u, int p) {
    tin_[u] = timer_++;                 // FIRST-IN: arriving at u
    for (int v : adj[u]) if (v != p)    // the parent guard (Ch.4, Trap 1)
        euler(v, u);
    tout_[u] = timer_++;                // LAST-OUT: backing out of u
}
// NOTE: this variant uses a 0..2n-1 clock. The flattened-array view
// (tin = preorder rank in 0..n-1, tout = tin + size) is Section 5.2.

bool isAncestor(int u, int v) {         // u an ancestor of v (u != v)?
    return tin_[u] <= tin_[v] && tout_[v] <= tout_[u];
}
\end{lstlisting}
\begin{complexitybox}
Time $\bigO{n}$ -- the clock is touched twice per node, the parent guard is checked once per
stored arc. Space $\bigO{n}$ for the three arrays plus $\bigO{h}$ recursion. On a
$2\cdot10^{5}$-node chain the recursion is $\bigO{n}$ deep and can RE -- run the same stamp
with an explicit stack (push \code{(node, parent, stage)}; \code{stage} $0$ = just entered,
$1$ = about to exit), or an iterative post-order. With a $0..n-1$ preorder clock and
$\tout=\tin+\size$ you get the exact \secref{c24-flatten} array for free.
\end{complexitybox}

\subsection{Worked Example}

\begin{example}{Stamping $G$ (children ascending)}{c24-times-run}
DFS from $0$ over $G$, visiting each node's neighbours in increasing order, on a clock that
ticks at entry and again at exit:
\[
  \begin{array}{c|ccccccc}
  \text{node} & 0 & 1 & 4 & 5 & 2 & 3 & 6\\ \hline
  \tin & 0 & 1 & 2 & 4 & 7 & 9 & 10\\
  \tout& 13& 6 & 3 & 5 & 8 & 12& 11
  \end{array}
\]
Order the rows by node id: $\tin=[0,1,7,9,2,4,10]$, $\tout=[13,6,8,12,3,5,11]$. Check
$\subtree(1)$: $1$ is entered at $1$, its children $4,5$ occupy $2..5$, $1$ exits at $6$ --
so $[\tin(1),\tout(1))=[1,6]$ covers exactly $\{1,4,5\}$. Ancestry: $0$ anc of $6$ since
$0\le 10$ and $11 \le 13$; $1$ \emph{not} anc of $2$ since $\tout(2)=8 > \tout(1)=6$. The
bracket string (open on entry, close on exit) is
$\code{( ( () () ) ( ) ( () ) )}$, i.e.\ \code{((()())(()())())}.
\end{example}
\begin{tracebox}[Reading the clock on the visit order $0,1,4,5,2,3,6$]
Enter $0$ ($\tin 0$). Enter $1$ ($\tin 1$). Enter $4$ ($\tin 2$), exit $4$ ($\tout 3$).
Enter $5$ ($\tin 4$), exit $5$ ($\tout 5$). Exit $1$ ($\tout 6$). Enter $2$ ($\tin 7$), exit
$2$ ($\tout 8$). Enter $3$ ($\tin 9$). Enter $6$ ($\tin{10}$), exit $6$ ($\tout{11}$), exit
$3$ ($\tout{12}$). Exit $0$ ($\tout{13}$). Max clock $=14=2n$, every pair nests or is
disjoint as \thmref{c24-anc} promises. \checkmark\ \figref{c24-intervals}.
\end{tracebox}

\subsection{Variations and Twists}

\begin{itemize}[nosep]
  \item \textbf{Preorder-rank clock} (this book's \secref{c24-flatten} choice): tick only on
        entry, $\tin\in[0,n)$, and set $\tout=\tin+\size$. Shorter intervals, directly index
        an array. \hfill\lv{2}
  \item \textbf{Both-tick clock} (shown above): $\tin,\tout\in[0,2n)$; the ``gap''
        $\tout-\tin$ still equals subtree span, handy for contour/perimeter reasoning. \hfill\lv{3}
  \item \textbf{1-indexed times}: judges sometimes want positions in $1..n$ -- add one
        everywhere; the containment inequalities are unchanged. \hfill\lv{2}
  \item \textbf{Bracket decoding}: given a balanced string, a stack rebuilds the tree -- the
        inverse map of \propref{c24-bracket} (\probref{c24-throne} leans on the forward map). \hfill\lv{3}--\lv{4}
  \item \textbf{Order-statistics on the walk}: sort nodes by $\tin$ to list them in the exact
        DFS encounter order, then filter (the shape of \probref{c24-throne}). \hfill\lv{3}
\end{itemize}

\subsection{Common Pitfalls}\label{sec:c24-times-pitfalls}

\begin{warningbox}[Trap 1 -- comparing only the entry time]
``$\tin(u) < \tin(v)$ so $u$ is above $v$'' is \emph{false}: a node entered earlier but in a
different branch is not an ancestor. You must also require $\tout(v) \le \tout(u)$ (or
$\tin(v) < \tout(u)$ in the preorder-rank form).
\end{warningbox}
\begin{lstlisting}[caption={\textcolor{warningred}{WRONG}: entry time alone is not containment}, language=C++, basicstyle=\ttfamily\small]
bool isAnc(int u, int v) { return tin_[u] < tin_[v]; }  // 1 is "anc" of 2 by this test -- it is not
\end{lstlisting}
\begin{lstlisting}[caption={\textcolor{successgreen}{RIGHT}: require the exit time too}, language=C++, basicstyle=\ttfamily\small]
bool isAnc(int u, int v) { return tin_[u] <= tin_[v] && tout_[v] <= tout_[u]; }
\end{lstlisting}

\begin{warningbox}[Trap 2 -- the exit clock never ticks (all $\tout$ wrong)]
Forgetting to increment on \emph{exit} leaves every $\tout$ at a stale value, so disjoint
subtrees look nested and the containment test fires on cousins.
\end{warningbox}
\begin{lstlisting}[caption={\textcolor{warningred}{WRONG}: stamps \code{tout} but never advances the clock}, language=C++, basicstyle=\ttfamily\small]
void euler(int u, int p) {
    tin_[u] = timer_++;
    for (int v : adj[u]) if (v != p) euler(v, u);
    tout_[u] = timer_;            // MISSING ++: no node's exit gets its own time
}
\end{lstlisting}
\begin{lstlisting}[caption={\textcolor{successgreen}{RIGHT}: advance the clock on the way out}, language=C++, basicstyle=\ttfamily\small]
void euler(int u, int p) {
    tin_[u] = timer_++;
    for (int v : adj[u]) if (v != p) euler(v, u);
    tout_[u] = timer_++;          // LAST-OUT tick, distinct from every entry
}
\end{lstlisting}

\begin{warningbox}[Trap 3 -- recursion depth on a chain-shaped tree]
Like rooting, the recursive stamp is $\bigO{h}$ deep; on a $2\cdot10^5$ path it overflows the
default stack and the run dies looking like a wrong answer. Do it iteratively.
\end{warningbox}
\begin{lstlisting}[caption={\textcolor{warningred}{WRONG}: recursion on a chain -> RE}, language=C++, basicstyle=\ttfamily\small]
euler(root, -1);   // depth ~ n on a path; blows the ~1e5-frame stack
\end{lstlisting}
\begin{lstlisting}[caption={\textcolor{successgreen}{RIGHT}: explicit stack with an entry/exit stage}, language=C++, basicstyle=\ttfamily\small]
// push {u, p, 0}; on pop with stage 0 record tin and push {u,p,1} then children {v,u,0};
// on pop with stage 1 record tout.  Same O(n), no recursion.
\end{lstlisting}

\subsection{Problem Lab}

One canonical problem anchors the section. LC 1600 \emph{Throne Inheritance} hands you a
family tree and asks for the line of succession -- which is \emph{exactly} the preorder
($\tin$ order) with the deceased filtered out. It is the ``order-statistics on the walk''
variation in problem clothing.

% ---------------------------------------------------------------------------
\solvedproblem{c24-throne}{Throne Inheritance}{LeetCode 1600}
  {https://leetcode.com/problems/throne-inheritance/}{LC Medium L3}
  {rooted tree, preorder = line of succession, filter-dead DFS}

\begin{problemstatement}[Throne Inheritance]
A king with name \code{kingName} founds a dynasty. Support three operations on the family
tree:
\begin{itemize}[nosep]
  \item \code{birth(parentName, childName)} -- \code{childName} is born as the
        \textbf{last} child of \code{parentName} (birth order matters);
  \item \code{death(personName)} -- marks that person dead (they stay in the tree);
  \item \code{getInheritanceOrder()} -- return the names, in order, of everyone still
        alive in the line of succession: the preorder (depth-first, children in birth order)
        traversal with dead people skipped.
\end{itemize}
No two people share a name; $1 \le$ total calls $\le 10^5$.
\end{problemstatement}

\begin{attemptgate}[Attempt it yourself first -- \lv{3}, give yourself 20 minutes]
A fair first attempt: model the family as an adjacency structure keyed by name, then think
about \emph{which} traversal the succession line is.
\peek{``eldest first, then recurse into each child before the next sibling'' -- which of the
four DFS orders from \chref{c03} is that?}
\nudge{store children in birth order; a single preorder DFS from the king lists the whole
line. Where do you drop the dead?}
\fullhint{Keep \code{children[name]} in insertion (birth) order and a
\code{dead} set. \code{getInheritanceOrder} $=$ preorder DFS from the king that appends a
name only if it is not in \code{dead}. Re-walk on each query (or cache) -- the order
\emph{is} the $\tin$ order. Full build below.}
\end{attemptgate}

\solstep{Understanding the Problem}
The succession order is a \term{preorder traversal} (``first-in'' order) of the family tree,
visiting children oldest-to-youngest, simply with deceased members omitted. This is the
$\tin$-order of \secref{c24-times} read as a list: node $u$ precedes node $v$ iff $\tin(u) <
\tin(v)$. So the whole problem is ``maintain an ordered children list under births, a dead
set under deaths, and dump the preorder of the living on demand.''

\solstep{Key Observations}
\begin{observation}
Births append children in order, so \code{children[parent]} is naturally a
\emph{birth-ordered} list -- iterating it left-to-right during a DFS yields the succession
order directly, no sorting needed.
\end{observation}
\begin{observation}
Death never removes a node from the tree (a dead person may have living descendants who are
still in line, after their elders) -- so keep the node and merely \emph{filter} it during the
walk.
\end{observation}
\begin{observation}
The succession list is exactly ``the preorder ranks $\tin$ sorted ascending, minus the dead'',
tying this problem to the timestamp machinery: there is nothing to recompute per query except
one DFS.
\end{observation}

\solstep{Think It Through With Me}
\begin{thinkbox}
\deadend{Re-derive the order with a priority queue keyed by age at every
\code{getInheritanceOrder}.}{The order is determined by \emph{tree structure + birth order},
not by absolute ages; a global sort ignores that a younger uncle precedes his own nephew only
if the uncle is an elder sibling of the nephew's parent -- a preorder, not an age sort.}
\deadend{On each birth/death, splice a flat global array of names.}{Correct-ish but you lose
the parent/child nesting that defines succession; keeping a \code{children} map per parent is
both simpler and $O(1)$ per birth.}
\aha{Store \code{children} in birth order and a \code{dead} hash set; \code{getInheritanceOrder}
is one preorder DFS from the king that skips dead names -- the $\tin$ order, filtered.}
\end{thinkbox}

\solstep{Solution}
Map each name to its ordered list of children; a set records the dead. On
\code{getInheritanceOrder}, DFS the king: append the current name if alive, then recurse into
children in stored (birth) order. Return the accumulated vector.

\solstep{Proof of Correctness}
\begin{proofbox}
By construction the children of each parent are held in birth order, so a DFS that visits
each node before its children (preorder) and each child before the next sibling emits nodes
in exactly the line-of-successence order defined by the problem (eldest branch exhausted
before the next). The DFS visits a node iff it is a descendant of the king, i.e.\ a dynasty
member; the algorithm appends precisely those visited nodes that are not in \code{dead},
matching ``everyone alive, in succession order''. Death keeps the node in the tree, so a dead
person's living descendants are still reached and appended at the right rank. Preorder visits
each node once (\thmref{c04-dfs-orient}), so nothing is duplicated or missed.
\end{proofbox}

\solstep{Implementation}
\begin{lstlisting}[caption={Throne Inheritance via ordered children + preorder DFS (LeetCode 1600)}, label={lst:c24-throne}]
#include <bits/stdc++.h>
using namespace std;

class ThroneInheritance {
public:
    ThroneInheritance(string kingName) : king_(kingName) {}

    void birth(string parentName, string childName) {
        children_[parentName].push_back(childName);   // append => birth order
    }
    void death(string name) { dead_.insert(name); }   // keep node, just mark

    vector<string> getInheritanceOrder() {
        vector<string> out;
        dfs(king_, out);
        return out;
    }

private:
    string king_;
    unordered_map<string, vector<string>> children_;
    unordered_set<string> dead_;

    void dfs(const string& u, vector<string>& out) {
        if (!dead_.count(u)) out.push_back(u);        // preorder: self first
        for (const string& c : children_[u])          // children in birth order
            dfs(c, out);
    }
};
\end{lstlisting}

\solstep{Trace}
\begin{tracebox}[King; king's children born $p_1, p_2$ in order; $p_1$ has child $c$; $p_2$ dead]
\code{children\_}: king$\to[p_1,p_2]$, $p_1\to[c]$. \code{dead\_}$=\{p_2\}$.
Preorder DFS from king: emit king; into $p_1$: emit $p_1$; into $c$: emit $c$; back; into
$p_2$: dead, skip $p_2$ (it has no children). Result $[\text{king}, p_1, c]$. This is the
$\tin$-ascending order with $p_2$ filtered -- exactly the timestamps of \secref{c24-times}
used as a list. \checkmark\ \probref{c24-throne}.
\end{tracebox}

\solstep{Complexity}
\begin{complexitybox}
\code{birth} and \code{death} are $\bigO{1}$ amortised (hash append / insert).
\code{getInheritanceOrder} is $\bigO{n}$ (one preorder pass, $n$ dynasty members). Space
$\bigO{n}$. If queries vastly outnumber births you can cache the vector and invalidate on
change, but the $\bigO{n}$ walk is accepted. The DFS is $\bigO{h}$ deep -- on a deep
lineage, convert to an explicit stack (Trap~3). Verdict: \ac.
\end{complexitybox}

\solstep{Pitfalls and Variants}
\begin{warningbox}
Using an \emph{unordered} container for children and then ``sorting by age'' instead of
relying on insertion order reintroduces the age-sort trap. \code{vector<string>} per parent
preserves birth order for free. And skipping a dead node must \emph{not} skip its subtree --
a dead parent has living heirs; only the name is filtered, the recursion continues.
\end{warningbox}

\begin{insightbox}[Key takeaways]
\begin{enumerate}[nosep, leftmargin=1.6em]
  \item ``Line of succession'' \emph{is} preorder $=\tin$ order -- recognise the traversal,
        and the problem is one DFS.
  \item Deletion that keeps descendants in the ordering (``filter, don't prune'') is a
        recurring pattern; mark and skip at emission time.
  \item Insertion-ordered children turn a would-be sort into a free invariant.
\end{enumerate}
\end{insightbox}

\begin{hardeningbox}
\begin{itemize}[nosep, leftmargin=1.4em]
  \item \textbf{Depth}: a $10^5$-deep single-line family makes the recursion overflow -- use
        an explicit stack, or reserve thread stack; the emit rule is unchanged.
  \item \textbf{String keys}: \code{unordered\_map} on \code{string} is fine; if the judge
        renames to integer ids, the same code with \code{int} keys is faster.
\end{itemize}
\end{hardeningbox}

\subsection{Practice Problems}

\begin{stuckbox}
\begin{itemize}[nosep, leftmargin=1.4em]
  \item Did you tick the clock on \emph{both} entry and exit (or use $\tout=\tin+\size$)?
        (Trap 2.)
  \item Does your ancestry test check $\tout$ as well as $\tin$? (Trap 1,
        \thmref{c24-anc}.)
  \item When filtering nodes (dead, banned), did you keep recursing into the survivors'
        children? (\probref{c24-throne}.)
  \item On a chain, is the stamp \emph{iterative} rather than $\bigO{n}$ recursion? (Trap 3.)
\end{itemize}
\end{stuckbox}

\begin{practiceset}
  \practice{Binary Tree Coloring Game}{LeetCode 1145}
    {https://leetcode.com/problems/binary-tree-coloring-game/}{LC Medium L3}
    {Two players colour nodes of a binary tree; the second wins iff some reachable region
    has more nodes than the rest. Find the value $\size[v]$ of the node the first player
    colours and compare it to $n-\size[v]$ -- a \emph{subtree size} question that the
    bottom-up $\size$ (and hence the $\tout-\tin$ span) answers directly. Root at the given
    value, DFS for its subtree size, and check whether any neighbour region exceeds the
    complement.}
\end{practiceset}

\begin{recallbox}[10-second refresher -- 5.1 Times \& Bracket Sequence]
One DFS stamps a clock at \emph{first-in} ($\tin$) and \emph{last-out} ($\tout$). Then $u$
is an ancestor of $v$ iff $\tin(u) \le \tin(v) < \tout(u)$ (\thmref{c24-anc}) -- $O(1)$.
Erase labels, print \code{(}/\code{)} and the tree is a balanced bracket string
(\propref{c24-bracket}); preorder $=\tin$ order (\probref{c24-throne}).
\selftest{
  \item Write the DFS that stamps $\tin$ and $\tout$ on a single clock from memory
        (\coderef{c24-stamp}).
  \item State and re-prove the containment test for ``$u$ is an ancestor of $v$''.
  \item On $G$ with children ascending, give $\tin$ and $\tout$, and say whether $1$ is an
        ancestor of $2$ (use both times).
  \item Write the bracket string for $G$'s rooting and read off which pairs match.
  \item Explain why a dead node in \probref{c24-throne} is skipped at emission but not pruned.
}
\end{recallbox}
